\documentclass[pdfbookmarks,a4paper,11pt]{article}
\usepackage{dsekcommon}
\usepackage[utf8]{inputenc}
\usepackage{enumerate}
\usepackage{lastpage}
\usepackage{longtable}
\ifpdf
  \usepackage{thumbpdf}
\fi
\usepackage{calc}
\usepackage[normalem]{ulem}
\usepackage{soul}
\usepackage{hyperref}
\hypersetup{
    colorlinks=true,
    linkcolor=black,
    urlcolor=blue
}

\pagestyle{fancy}

\lhead{}
\chead{}
\rhead{}
\lfoot{}
\cfoot{\thepage\ (\nohyperpageref{LastPage})}
\rfoot{}

\usepackage[T1]{fontenc}
\usepackage[utf8]{inputenc}
\usepackage[swedish]{babel}
\usepackage{multicol}
\usepackage{xcolor}
\usepackage{tabularx}
\usepackage{booktabs}

\setheader{Krishanteringsplan}{Riktlinje}{}

\title{Krishanteringsplan}

\begin{document}

\maketitle

\tableofcontents

\newpage

\section{Formalia}

\subsection{Sammanfattning}

Denna krishanteringsplan beskriver hur sektionen bör agera i det fall att en kris uppstår. Planen ska fungera som ett stöd för sektionens medlemmar och innehåller riktlinjer för åtgärder vid olika typer av krissituationer. Målet är att säkerställa ett snabbt, strukturerat och effektivt agerande för att minimera skador och upprätthålla trygghet inom verksamheten. 

\subsection{Syfte}

Syftet med denna krishanteringsplan är att skapa en gemensam struktur och tydliga rutiner för hur sektionen ska hantera och följa upp kriser. Planen ska säkerställa att medlemmar och funktionärer kan fatta välgrundade beslut och agera samordnat, samt bidra till att minska osäkerhet och förebygga att kriser eskalerar eller får onödigt stora konsekvenser. Sektionens krishanteringsplan fungerar som ett komplement till Teknologkårens krishanteringsplan, som finns att hitta bland \href{https://www.tlth.se/styrdokument}{Teknologkårens styrdokument}.

\subsection{Omfattning}

Riktlinjen omfattar sektionen i sin helhet. 

\subsection{Historik}

Skriven av Ella Svensson 2026. 

\section{Allmänt om kriser}

\textcolor{red}{\textbf{Vid akuta nödsituationer - ring alltid 112.}}

Ingen funktionär förväntas att vara livräddare, men att känna till hur man kan agera i en krissituation kan göra stor skillnad för alla berörda. Tveka aldrig att kontakta professionella aktörer, såsom sjukvårdspersonal, polis, räddningstjänst och väktare för stöd och vägledning. Utsätt aldrig dig själv eller andra för onödiga risker i ett försök att hjälpa någon annan. Tänk alltid på din egen säkerhet och mående när krisen uppstår!

\textbf{Vid alla krissituationer, stora som små, ska sektionsordförande och kårordförande informeras så snart situationen tillåter:}

\begin{itemize}
    \item Sektionsordförande Ella Svensson, 070-304 68 92 eller ordforande@dsek.se
    \item Kårordförande Axel Kjellberg, 070-418 79 21 eller ko@tlth.se
\end{itemize}

\section{Allmän handlingsplan vid krissituationer}

Följande avsnitt beskriver en generell handlingsplan som stöd vid hantering av en krissituation. Eftersom varje situation är unik kan vissa åtgärder behöva anpassas utifrån omständigheterna. För hantering av specifika händelser, som brand, diskriminering, olycka eller dödsfall, se avsnitt 4.

\subsection{Bedöm situationen och skapa trygghet}

Det första steget är att skapa en lägesbild och säkerställa att ingen befinner sig i omedelbar fara. Försök agera, inte reagera!

\begin{itemize}
    \item Få en överblick av situationen - vad? Var? Vem?
    \item Kontrollera om någon är skadad eller behöver akut hjälp
    \item Säkerställ din egen och andras säkerhet och avlägsna personer från fara om det kan göras utan risk
    \item Skapa lugn och trygghet för de berörda
\end{itemize}

\subsection{Larma och kontakta hjälp}

Vid behov ska professionella aktörer kontaktas så tidigt som möjligt.

\begin{itemize}
    \item Ring \textbf{112} vid fara för liv, allvarlig skada eller andra nödsituationer
    \item Följ instruktionerna från larmoperatören
    \item Hämta första hjälpen-utrustning och hjärtstartare om situationen kräver det och dessa finns tillgängliga
    \item Kontakta och ta hjälp av vakter eller annan relevant personal vid behov
\end{itemize}


\subsection{Informera, samordna och stötta}

När situationen är under kontroll ska relevanta personer informeras och den fortsatta hanteringen samordnas. Fokus ska även ligga på att stötta de personer som har påverkats av händelsen.

\begin{itemize}
    \item Kontakta sektionsordförande och kårordförande, kontaktuppgifter finns under avsnitt 7!
    \item Informera andra relevanta personer utifrån situationens omfattning och fördela ansvar
    \item Ta hand om och stötta personer som har bevittnat eller påverkats av händelsen. Se avsnitt 5 för de stödresurser som finns!
    \item Visa hänsyn till de berördas integritet och undvik att dela känslig information eller spekulationer
\end{itemize}

\subsection{Efterarbete och uppföljning}

Efter att den akuta situationen är hanterad är det viktigt att följa upp händelsen, ge fortsatt stöd till berörda personer och ta tillvara på erfarenheter inför framtiden.

\begin{itemize}
    \item Säkerställ att berörda personer får möjlighet till fortsatt stöd
    \item Dokumentera händelseförloppet, fattade beslut och vidtagna åtgärder
    \item Säkerställ att relevant information har lämnats vidare till ansvariga personer
    \item Sektionsstyrelsen ansvarar för att utvärdera krishanteringen och vid behov uppdatera rutiner
\end{itemize}

\section{Hantering av specifika händelser}

\subsection{Brand}

En brand kan snabbt utvecklas och det är därför viktigt att agera direkt. Vid brand ska allas säkerhet alltid prioriteras framför att försöka släcka elden eller rädda material. Använd principen Rädda – Varna – Larma – Släck som vägledning och följ stegen nedan: 

\begin{enumerate}
    \item Rädda personer som befinner sig i omedelbar fara, om det kan göras utan att utsätta dig själv eller andra för risk
    \item Varna personer i närheten och påbörja utrymning av lokalen
    \item Aktivera brandlarmet om det inte redan har utlösts
    \item Ring 112 och uppge vad som har hänt samt var branden befinner sig. Alla adresser hittas i tabellen nedan!
    \item Om branden är liten och kan släckas utan risk för personskada kan brandsläckare eller annan släckutrustning användas
    \item Informera räddningstjänsten om det finns personer som kan ha svårt att utrymma byggnaden eller som misstänks vara kvar
    \item Samlas på återsamlingsplatsen och kontrollera, om möjligt, att samtliga personer har tagit sig ut och är i säkerhet
    \item Kontakta sektionsordförande och kårordförande och informera om händelsen
\end{enumerate}

\textbf{Kom ihåg:} Att brandlarmet tystnar innebär inte att faran är över, räddningstjänsten kan stänga av larmet under sitt arbete. Återvänd aldrig till byggnaden förrän räddningstjänsten har meddelat att det är säkert!

\begin{table}[h]
\centering
\renewcommand{\arraystretch}{1.3} % Lite högre rader

\begin{tabularx}{\textwidth}{l X X}
\toprule
\textbf{Plats} & \textbf{Adress} & \textbf{Återsamlingsplats} \\
\midrule
E-Huset & Kommunikationsvägen 2, 223 63 & Parkeringen på baksidan \\
Lophtet & Getingevägen 66, 222 41 & Parkeringen mitt emot KC:4 \\
Kårhuset & John Ericssons väg 3, 223 63 & Navet \\
\bottomrule
\end{tabularx}
\end{table}

\subsection{Olycka eller sjukdom}

En olycka eller ett akut sjukdomsfall kan inträffa plötsligt och kräver ofta ett snabbt och lugnt agerande. Fokusera på att skapa en säker situation, tillkalla hjälp och ge stöd till den drabbade, försök inte ge vård utöver din kunskap eller förmåga! Om du befinner dig på plats bör du göra följande:

\begin{enumerate}
    \item Säkerställ att platsen är säker för både dig själv och den drabbade
    \item Bedöm personens tillstånd och kontrollera om personen är vid medvetande och andas normalt
    \item Ring 112 vid misstanke om allvarlig sjukdom eller skada, eller om du är osäker på hur allvarlig situationen är. Följ larmoperatörens instruktioner!
    \item Ge första hjälpen utifrån din kunskap och utbildning. Be någon hämta första hjälpen-utrustning eller hjärtstartare om situationen kräver det och dessa finns tillgängliga
    \item Skapa lugn kring den drabbade och ge utrymme för eventuell räddningspersonal att arbeta
    \item Informera ansvariga funktionärer, vakter eller annan relevant personal om situationen
    \item Om personen förs till sjukhus, säkerställ om möjligt att personen har sällskap av en vän, anhörig eller annan lämplig person
    \item Kontakta sektionsordförande och kårordförande och informera om händelsen
\end{enumerate}

\textbf{Kom ihåg:} Lämna inte den drabbade ensam om det inte är nödvändigt för att hämta hjälp eller säkerställa din egen säkerhet! Om situationen tillåter kan det även vara bra att notera den drabbades kontaktuppgifter för den fortsatta hanteringen.

\subsection{Trakasserier och dikriminering}

Trakasserier och diskriminering kan ta sig många olika uttryck och kan förekomma både i fysiska och digitala sammanhang. Om du misstänker att någon har blivit utsatt ska situationen alltid tas på allvar, den utsattas trygghet och integritet ska alltid stå i fokus! Vid misstanke om trakasserier eller diskriminering gäller följande:

\begin{enumerate}
    \item Säkerställ att den utsatta personen befinner sig i en trygg miljö. Om situationen pågår, avbryt den om det kan göras utan att utsätta dig själv eller andra för fara
    \item Kontakta 112 vid akuta situationer eller om någon är i omedelbar fara. Vid arrangemang där vakter finns ska dessa också kontaktas
    \item Ta hand om den utsatta personen och lyssna på vad personen känner att den behöver. Se till att personen kommer hem på ett säkert sätt och har sällskap. Teknologkåren bekostar taxiresa om behov finns
    \item Dokumentera relevanta uppgifter om händelsen när situationen tillåter, exempelvis signalement eller kontaktuppgifter. Spara även skärmdumpar om trakasserierna har skett digitalt!
    \item Kontakta sektionsordförande och kårordförande och informera om händelsen
\end{enumerate}

\textbf{Kom ihåg:} Visa respekt för den utsattas integritet! Sprid inte information om händelsen till andra än de personer som behöver känna till den för att hantera situationen. Om du själv påverkas av händelsen eller hanteringen av den finns det alltid stöd att få genom universitetets stödresurser eller sektionens Trivselråd!

\subsection{Våld och hot}

Hotfulla eller våldsamma situationer kan snabbt förändras och det är viktigt att alltid prioritera den egna säkerheten. Ingen funktionär förväntas utsätta sig själv för fara genom att ingripa fysiskt. I många situationer är det bästa agerandet att skapa avstånd, tillkalla hjälp och använda de resurser som finns tillgängliga. Befinner du dig på plats bör du göra följande:

\begin{enumerate}
    \item Säkerställ din egen och andras säkerhet. Undvik att ingripa fysiskt om det innebär en risk för dig själv eller andra
    \item Ring alltid 112 vid akuta eller allvarliga situationer. Uppge var ni befinner er och följ instruktionerna från larmoperatören
    \item Kontakta vakter om dessa finns tillgängliga
    \item Om situationen tillåter, hjälp den utsatta personen bort från situationen och till en trygg plats
    \item Ta hand om den utsatta personen och säkerställ att personen får den hjälp som behövs, exempelvis sjukvård. Se till att personen kommer hem på ett säkert sätt och har sällskap, Teknologkåren bekostar taxiresa om behov finns
    \item Dokumentera relevanta uppgifter om händelsen när situationen tillåter, exempelvis signalement eller kontaktuppgifter
    \item Vid digitala hot eller trakasserier, spara relevanta bevis såsom skärmdumpar
    \item Kontakta sektionsordförande och kårordförande och informera om händelsen
\end{enumerate}

\textbf{Kom ihåg:} Visa respekt för den utsattas integritet! Sprid inte information om händelsen till andra än de personer som behöver känna till den för att hantera situationen. Om du själv påverkas av händelsen eller hanteringen av den finns det alltid stöd att få genom universitetets stödresurser eller sektionens Trivselråd!

\subsection{Dödsfall}

Ett dödsfall är en mycket allvarlig händelse som kräver ett lugnt och strukturerat agerande. Om en person inte går att rädda ska fokus ligga på att ta hand om närvarande personer, säkerställa att rätt personer kontaktas och skapa trygghet i situationen. Vid dödsfall har Teknologkåren en övergripande krishanteringsplan och kan bistå sektionen med stöd och resurser. Vid ett misstänkt eller bekräftat dödsfall bör följande åtgärder vidtas:

\begin{enumerate}
    \item Kontakta 112 vid upptäckt av ett misstänkt dödsfall eller om en person är livlös. Följ instruktionerna från larmoperatören
    \item Kontakta omedelbart kårordförande och sektionsordförande. Vid behov tar kåren över eller stödjer den fortsatta krishanteringen
    \item Skapa lugn och trygghet för personer som har bevittnat händelsen eller påverkats av den
    \item Samla, om möjligt och när situationen tillåter, kontaktuppgifter till personer som kan behöva fortsatt stöd eller information
    \item Låt platsen vara orörd i väntan på instruktioner från polis eller annan ansvarig personal
    \item Sprid inte information om händelsen innan ansvariga personer har samordnat kommunikationen
\end{enumerate}

\textbf{Kom ihåg:} Att bevittna eller hantera ett dödsfall kan vara en mycket svår upplevelse. Även personer som varit på plats och försökt hjälpa till kan påverkas av händelsen. Tveka inte att söka stöd vid behov, se avsnitt 5 för tillgängliga stödresurser!

\section{Stödresurser}

Efter en krissituation kan både direkt berörda personer och personer som varit involverade i hanteringen behöva stöd. Det är viktigt att komma ihåg att även den som agerat som funktionär eller försökt hjälpa till kan påverkas av en händelse. Nedan finns information om tillgängliga stödresurser inom och utanför sektionen.

\subsection{Trivselrådet}

Trivselrådet är en stödresurs inom sektionen som arbetar med frågor kopplade till likabehandling, internationalisering och arbetsmiljö. Trivselrådet finns till för alla sektionsmedlemmar och kan kontaktas om något har hänt, om man känner sig utsatt eller om man behöver någon att prata med. 

Vid krissituationer som exempelvis rör trakasserier, diskriminering, konflikter eller andra händelser som påverkar en medlems trygghet kan Trivselrådet bidra med stöd och vägledning. Ärenden och funderingar kan lämnas in via Trivselrådets formulär, där det även finns möjlighet att vara anonym. Både Likabehandlingsombuden och Trivselmästaren omfattas av tystnadsplikt.

Likabehandlingsombuden kan kontaktas via likabehandlingsombud@dsek.se, och dessa mejl tas inte del av Trivselmästaren om ärendet skulle gälla denne. Trivselmästaren kan kontaktas via trivselm@dsek.se.

\textbf{Trivselrådets kontaktformulär:} \href{https://bit.ly/trivselkontakt}{bit.ly/trivselkontakt}

\subsection{Teknologkåren}

Teknologkåren är studentkårens gemensamma stödorganisation och finns till för att stötta sektionerna och deras medlemmar. Vid större eller mer omfattande krissituationer kan Teknologkåren bidra med stöd, rådgivning och samordning i den fortsatta hanteringen. Kåren kan även hjälpa till vid situationer som kräver extra stöd kring kommunikation, kontakt med universitetet eller andra externa aktörer. Vid allvarliga händelser ska kårordförande alltid kontaktas för att säkerställa att rätt resurser och stödinsatser kan aktiveras. 

\textbf{Mer information och kontaktuppgifter:} \href{https://www.tlth.se/}{Teknologkåren}

\subsection{Studenthälsan}

Studenthälsan erbjuder professionellt stöd till studenter vid Lunds universitet som behöver hjälp med sitt mående eller upplever psykiska påfrestningar. Den som har påverkats av en krissituation, antingen direkt eller genom att ha bevittnat en händelse, kan vända sig till Studenthälsan för stöd och rådgivning. Studenthälsan erbjuder bland annat individuella samtal och all kontakt är kostnadsfri samt omfattas av tystnadsplikt. 

\textbf{Mer information och kontaktuppgifter:} \href{https://www.lu.se/student/studentstod-och-halsa/studenthalsan}{Studenthälsan | Lunds universitet}

\subsection{Studentprästerna}

Studentprästerna erbjuder samtalsstöd till studenter vid Lunds universitet och finns till för alla, oavsett tro eller livsåskådning. De erbjuder enskilda samtal för den som vill prata om exempelvis kris, sorg, oro, relationer eller andra frågor som känns viktiga. Samtalen är kostnadsfria och Studentprästerna har tystnadsplikt.

\textbf{Mer information och kontaktuppgifter:} \href{https://www.svenskakyrkan.se/lund/studentprasterna}{Studentpräster vid Lunds universitet/Student chaplaincy at Lund University - Svenska kyrkan i Lund}

\section{Kommunikation och media}

Vid en krissituation är det viktigt att information hanteras strukturerat och med hänsyn till de berörda personernas integritet. För att undvika ryktesspridning och spridning av felaktig information ska den officiella kommunikationen om händelsen hanteras av utsedda personer.

\subsection{Intern kommunikation}

Sektionsordförande ansvarar för den interna kommunikationen till sektionens medlemmar och samordnar vilken information som ska delas, när den ska delas och på vilket sätt. Funktionärer och andra personer som varit involverade i hanteringen av en kris ska:

\begin{itemize}
    \item Inte sprida information om händelsen vidare utan godkännande från ansvariga personer
    \item Undvika spekulationer och endast dela bekräftad information
    \item Visa hänsyn till berörda personer och deras integritet
\end{itemize}

Vid kommunikation efter en krissituation ska fokus ligga på att ge korrekt information, skapa trygghet och informera om eventuella stödresurser.

\subsection{Extern kommunikation och media}

Vid en krissituation kan det finnas behov av att kommunicera med externa aktörer, exempelvis Teknologkåren, universitetet, media eller andra organisationer. För att säkerställa att informationen är korrekt och samordnad ansvarar utsedda personer för sektionens officiella externa kommunikation.

Kårordförande ansvarar för extern kommunikation och är sektionens kontaktperson gentemot media och andra externa aktörer vid större krissituationer. Kårordförande har erfarenhet av mediekontakter och kan besvara fler frågor än någon på sektionsnivå.

Funktionärer och medlemmar har rätt att uttrycka sina egna åsikter och erfarenheter, men ska vara tydliga med att dessa är personliga och inte representerar sektionens eller Teknologkårens officiella ståndpunkt. Vid kontakt med media rekommenderas att hänvisa vidare till kårordförande för frågor som rör sektionens eller kårens hantering av händelsen.

\section{Kontaktuppgifter}

\subsection{Akuta kontaktvägar}

\begin{table}[h]
\centering
\renewcommand{\arraystretch}{1.3}

\begin{tabularx}{\textwidth}{X l}
\toprule
\textbf{Funktion} & \textbf{Telefon} \\
\midrule
SOS Alarm & 112 \\
Polisen (ej akut) & 114 14 \\
Sjukvårdsupplysning & 1177 \\
LU:s larmtelefon & 046-222 07 00 \\
\bottomrule
\end{tabularx}
\end{table}

Genom LU:s larmtelefon kan du komma i kontakt med väktare, kuratorer och universitetets krishantering. Numret finns även på baksidan av alla LU-kort.

\subsection{Sektionen}

Kontaktuppgifterna uppdateras vid nyval av posterna.

\begin{table}[h]
\centering
\renewcommand{\arraystretch}{1.3}

\begin{tabularx}{\textwidth}{X l l X}
\toprule
\textbf{Roll} & \textbf{Namn} & \textbf{Telefon} & \textbf{Mail} \\
\midrule
Sektionsordförande & Ella Svensson & 070-304 68 92 & ordforande@dsek.se \\
Vice sektionsordförande & Loke Bahr & 072-157 79 95 & vordforande@dsek.se \\
Trivselmästare & Theo Ramsö & 072-536 02 80 & trivselm@dsek.se \\
\bottomrule
\end{tabularx}
\end{table}

\subsection{Teknologåren}

\begin{table}[h]
\centering
\renewcommand{\arraystretch}{1.3}

\begin{tabularx}{\textwidth}{X X X}
\toprule
\textbf{Roll} & \textbf{Telefon} & \textbf{Mail} \\
\midrule
Kårordförande & 070-418 79 21 & ko@tlth.se \\
Pedell & 070-418 79 29 & pedell@tlth.se \\
Aktivitetssamordnare & 070-418 79 30 & aktsam@tlth.se \\
Studiesocialt ansvarig & 070-418 79 25 & sa@tlth.se \\
Expeditionen & 070-418 79 20 & expen@tlth.se \\
\bottomrule
\end{tabularx}
\end{table}

\subsection{Övriga kontakter}

\begin{table}[h]
\centering
\renewcommand{\arraystretch}{1.3}

\begin{tabularx}{\textwidth}{l l X l}
\toprule
\textbf{Roll} & \textbf{Namn} & \textbf{Mail} & \textbf{Telefon} \\
\midrule
Husprefekt E-huset & Stefan Höst & \url{husprefekt@ehuset.lth.se} & 046-222 75 32 \\
Husintendent E-huset & Per-Henrik Rasmussen & \url{ph@ehuset.lth.se} & 046-222 74 93 \\
Kurator LTH & Matilda Andersson & \url{matilda.andersson@lth.lu.se} & 046-222 85 16 \\
Studentprästerna & -- & \url{studentprasterna.lund@svenskakyrkan.se} & 046-71 87 35 \\
Studenthälsan & -- & -- & 046-222 43 77 \\
\bottomrule
\end{tabularx}
\end{table}


\end{document}